%% -------------------------------------------------<< \readme 为 MUST Thesis 模版说明,正式使用时注釋掉
\newgeometry{left=0cm,right=0cm,top=0cm, bottom=0cm}
\thispagestyle{empty}
\vfill
\begin{center}
	\begin{tikzpicture}
		% 子图全局透明设置
		\begin{scope}[opacity=0.0]
			\draw[help lines,gray!20,step=1cm] (-0.49\paperwidth,-0.49\paperheight) grid (0.49\paperwidth,0.49\paperheight);
			\fill (0,0) circle (1pt);
			\foreach \i in {-10,-9,...,10}
			\node at (\i,0) {{\large \i}};
			\foreach \i in {-14,-13,...,14}
			\node at (0,\i) {{\large \i}};
		\end{scope}

		% GitHub 下载
		\coordinate (cd1) at (-7.6,13.5);
		\node[black!80](githubico) at (cd1) {\sfz{20pt}\faGithub};
		\node[right=-2.87pt of githubico.east,fill=black!80,rectangle, rounded corners=2pt] {\url{https://github.com/iihciyekub/mustthesis}};
		\draw[white,line width=1.3pt]  (cd1) circle (10pt);



		% 校徽logo
		\path
		let \p1 =(-4,2) in
		% 节点放置校徽图片
		node(logo)[anchor=east] at (\p1) {\includegraphics[width=3.5cm]{figure/MUSTSchoolBadgecolor}}
		%  澳门科技大学
		node(c)[anchor=west,text width =10cm,right=1pt of logo.15]  {\sfz{35pt}澳\hfill 門\hfill 科\hfill 技\hfill 大\hfill 學}


		% 学校英文名
		node(must)[anchor=west,text width=10cm,right=1pt of logo.-15] {\color{gray}\sfz{10pt}\textbf{MACAU \hfill UNIVERSITY\hfill OF\hfill SCIENCE\hfill AND\hfill TECHNOLOGY}}
		% 线段
		node[anchor= west,right=1pt of logo.-21]{\rule{10cm}{0.7pt}}
		node[anchor= west,right=1pt of logo.-22]{\rule{2.87cm}{1.6pt}}
		% 学院
		node(SOB)[anchor=north west,right=1pt of logo.-28,gray]{{\footnotesize School of Business,2018}}
		% 版本说明
		node[left=0pt of c.east,yshift=10mm]{\color{Salmon}\footnotesize\projectBname \color{black!60} \sfz{18pt}{\faGraduationCap} }


		node(L1)[anchor= north west,white,text height=5pt,fill=black!80,rectangle, rounded corners=1pt,right=2pt of logo.-37]
			{\href{https://www.latex-project.org/lppl/}{ \color{white}\sfz{7pt}\texttt{\rlap{\color{black!80}p}Licenes:}}}
		node(L2)[anchor= north west,text height=5pt,right = 0mm of L1,rectangle, fill=Periwinkle!80,rounded corners=1pt,right=1pt of L1] {\href{https://www.latex-project.org/lppl.txt}{\color{white}\sfz{7pt}\texttt{The LaTeX Project Public License (LPPL v1.3c)}}}
		node[anchor=north west,right=0pt of L2]{\faBalanceScale\faCreativeCommons}
		;





		% author 作者信息	
		\begin{scope}
			\tikzstyle{a}=[draw=black!40,line width=1.5pt]
			\coordinate (P1) at (5,-11);
			\clip (P1) circle (0.55cm);
			\node(note2) at (P1) {\includegraphics[width=1.2cm]{figure/author}};
			\draw[a]  (P1) circle (0.55cm);
		\end{scope}


		\node(e)[right=-1mm of note2.east,yshift=2mm] {};
		\node(ee2)[right=-1mm of note2.east,yshift=-12mm] {};
		\draw[Goldenrod,line width=3pt,name=l1] (e)--(ee2);

		\node(name1)[right=-.5mm of e ,anchor=north west] {\href{mailto:lyj8512@126.com}{\texttt{@}\iiname}};
		\node[right=-.5mm of name1.west,yshift=-1.2em,gray] {\color{LimeGreen}{\faWechat} \color{gray}\texttt{iichieykub}};
		\node[right=-.5mm of name1.west,yshift=-2.4em,gray] {{\footnotesize \texttt{Copyright} \copyright  \texttt{2019}}\iiname};

	\end{tikzpicture}
\end{center}

\vfill


%% 第二页
\newpage
\thispagestyle{empty}

\vfill
\begin{center}
	\begin{tikzpicture}

		%% -------------------------------------------------<< 子图全局透明设置
		\begin{scope}[opacity=0]
			\draw[help lines,gray!20,step=1cm] (-0.49\paperwidth,-0.49\paperheight) grid (0.49\paperwidth,0.49\paperheight);
			\fill (0,0) circle (1pt);
			\foreach \i in {-10,-9,...,10}
			\node at (\i,0) {{\large \i}};
			\foreach \i in {-14,-13,...,14}
			\node at (0,\i) {{\large \i}};
		\end{scope}




		%% -------------------------------------------------<<
		% 设置版本节点的样式
		\tikzset{gitnode/.style={anchor=east,fill=#1!50, draw=#1, text=black,rounded corners=2pt,opacity=0.7,node font=\texttt,font=\large,line width=0.6pt}}
		\tikzset{gitsubnode/.style={anchor=east,fill=#1!50, circle=1pt,draw=#1, text=black!60,line width=0.6pt}}

		% 设置全局坐标点
		\coordinate (git1) at (-5.5,12.5);
		\filldraw[gray!10!blue!2!red!1,opacity=0.2,rounded corners=5pt,draw=black,line width=.8pt,name path=A]
		let \p1=(git1) in
		(\x1-4cm,\y1+1.1cm) rectangle +(19.2cm,-5.78) ;

		% 发布版本节点的代码
		\node(m1)[gitnode=LimeGreen] at (git1){\texttt{Release}};
		\draw[dashed,line width=.8pt,gray!40]   ($(m1.east)+(3mm,0)$) -- +(2cm,0)node[gitsubnode=LimeGreen,name=mm1]{}
		--+(4.7cm,0)node[gitsubnode=LimeGreen,name=v2]{}
		--+(12.4cm,0)node[gitsubnode=LimeGreen,name=v3]{}
		+(13,0);

		% v0.0.0.1
		\node[anchor=east,right=-6pt of mm1.west,yshift=5mm]  at (mm1) {\faHandODown\texttt{~v0.0.0.1}};
		% v0.0.0.4
		\node[anchor=east,right=-6pt of v2.west,yshift=5mm]  at (v2) {\faHandODown\texttt{~v0.0.0.4}};
		% v0.0.1.8
		\node[anchor=east,right=-6pt of v3.west,yshift=5mm]  at (v3) {\faHandODown\texttt{~v0.0.1.8}};



		% Prof.Jenny
		\node(m2)[gitnode=Melon,yshift=-1cm,label=left:\texttt{Init}] at (git1){\texttt{@Prof.Jenny}};
		\draw[dashed,line width=.8pt,gray!40]   ($(m2.east)+(6mm,0)$) --+(.3cm,0)node[gitsubnode=Melon,name=mm2]{\sfz{7pt}\faFlash} --+(13,0);
		\draw[->>,Melon,line width=1.5pt](mm2.north) .. controls +(0,.8) and +(-.6,0) .. (mm1.west);



		% FP
		\node(m3)[gitnode=Rhodamine,yshift=-2cm,label=left:\texttt{deve}] at (git1){\texttt{@pīng.Fán}};
		\draw[dashed,line width=.8pt,gray!40]   ($(m3.east)+(3mm,0)$) --+(2cm,0)node[gitsubnode=Rhodamine,name=mm3]{\sfz{8pt}\faPagelines} --+(13,0);
		\draw[->,gray,line width=1.5pt,opacity=0.7](mm2)-- (mm3.north west);






		% YLF
		\node(m4)[gitnode=Turquoise,yshift=-3cm,label=left:\texttt{deve}] at (git1){\texttt{\texttt{@yàn.lì.Fāng}}};
		\draw[dashed,line width=.8pt,gray!40]   ($(m4.east)+(3mm,0)$) --+(3cm,0)node[gitsubnode=Turquoise,name=mm4]{\faPlug } --+(13,0);
		\draw[->,gray,line width=1.5pt,opacity=0.7](mm3)-- (mm4.north west);
		\draw[->>,Turquoise,line width=1.5pt](mm4.north) .. controls +(0,2.3) and +(-1,0) .. (v2.west);




		\node(m5)[gitnode=SeaGreen,yshift=-4cm,label=left:\texttt{deve}] at (git1){\texttt{@}\iiname};
		\draw[dashed,line width=.8pt,gray!40]   ($(m5.east)+(3mm,0)$) --+(10cm,0)node[gitsubnode=SeaGreen,name=mm5]{\faBug} --+(13,0);
		\draw[->,gray,line width=1.5pt,opacity=0.7](mm4)-- (mm5.west);
		\draw[->>,SeaGreen,line width=1.5pt](mm5.north) .. controls +(0,3.1) and +(-1,0) .. (v3.west);








		%% -------------------------------------------------<< 文件夹
		\coordinate (f2) at ($(git1)+(2,-6)$);
		% 文件夹背景图
		\filldraw[gray!10!blue!2!red!1,opacity=0.6,rounded corners=5pt,draw=black,line width=.8pt,name path=A]
		let \p1=(f2) in
		(\x1-6cm,\y1+1.1cm) rectangle +(13.2cm,-9.8) ;

		% 项目文件根目录
		\node(project)[label={above:\faHandORight ~~\color{Purple}\projectSname}] at (f2){{\huge \faFolderOpen}};


		% content
		\def\ff{folder1}
		\node(\ff)[
			right=0pt of project.south east, yshift=-.25cm,
			label={right:\color{blue}\texttt{\textbf{content}}}
		]
		{{\Large \faFolderOpenO}};


		\draw[line width=1pt,rounded corners=2pt,gray](project) |-(\ff) ;
		\path [name path=C](\ff) -- +(20,0);
		\path [name intersections={of=A and C,by=F}];
		\fill[name=ii,opacity=0.6,rounded corners=1pt] ($(F) + (-4pt,-2pt)$) rectangle +(8pt,4pt) node(ee){};
		\draw (ee.center) +(0,-2pt) edge[in=180,out=0] +(2,2)  edge[in=180,out=0] +(1,2);

		\foreach \a /\b /\c  in {
				25/\texttt{En-Abstract.tex}/\faFileTextO,
				75/\texttt{Zh-Abstract.tex}/\faFileTextO,
				125/\texttt{chapter01.tex}/\faFileTextO,
				175/\texttt{chapter02.tex}/\faFileTextO,
				225/\texttt{chapter03.tex}/\faFileTextO,
				275/\texttt{appendix.tex}/\faFileTextO,
				325/\texttt{acknowledgement.tex}/\faFileTextO
			}
			{
				\node(file\a)[right=0pt of \ff.south east,yshift=- \a * .1mm,label=right:\color{gray}\texttt{{\footnotesize \b}},black!70]{ \c };
				\draw[line width=1pt,rounded corners=2.25pt,gray](\ff) |-(file\a);
				%	\draw[line width =.6pt,gray](file\a.west) .. controls +(180:2cm) and +(0:3.1cm) .. +(-2.854cm, -25+\a * .1mm);
			}



		% figure
		\ff{folder2}
		\node(\ff)[left=0pt of project.south west, yshift=-.6cm,label={left:\color{blue}\texttt{\textbf{figure}}}]{{\Large \faFolderOpenO}};
		\draw[line width=1pt,rounded corners=2pt,gray](project) |-(\ff);

		\foreach \a /\b /\c  in {
				25/\texttt{MUSTSchoolBadgebk.pdf}/\faFilePdfO,
				75/\texttt{MUSTSchoolBadgecolor.pdf}/\faFilePdfO,
				125/\texttt{eg05.png}/\faFilePdfO,
				175/\texttt{eg06.png}/\faFileImageO,
				225/\texttt{author.jpg}/\faFileImageO
			}
			{
				\node(file\a)[left=0pt of \ff.south west,yshift=-\a * .1mm,label=left:\color{gray}\texttt{{\footnotesize \b}},black!70]{ \c };
				\draw[line width=1pt,rounded corners=2.25pt,gray](\ff) |-(file\a);
			}


		% fonts
		\ff{folder3}
		\node(\ff)[right=0pt of project.south east, yshift=-4.5cm,label={right:\color{blue}\texttt{\textbf{fonts}}}]{{\Large \faFolderOpenO}};
		\draw[line width=1pt,rounded corners=1pt,gray](project) |-(\ff);

		\foreach \a /\b /\c  in {
				25/\texttt{FontAwesome.otf}/\faFont,
				75/\texttt{KaiTi.ttf}/\faFont
				%		125/\texttt{MUSTThesistemplatever.pdf}/\faFilePdfO,	
			}
			{
				\node(file\a)[right=0pt of \ff.south east,yshift=-\a * .1mm,label=right:\color{gray}\texttt{{\footnotesize \b}},black!70]{ \c };
				\draw[line width=1pt,rounded corners=2.25pt,gray](\ff) |-(file\a);
			}




		% format
		\ff{folder4}
		\node(\ff)[left=0pt of project.south west, yshift=-3.75cm,label={left:\color{blue}\texttt{\textbf{format}}}]{{\Large \faFolderOpenO}};
		\draw[line width=1pt,rounded corners=2pt,gray](project) |-(\ff);

		\foreach \a /\b /\c  in {
				25/\texttt{fontawesome.sty}/\faFileCodeO,
				75/\texttt{MUSTThesisC.cls}/\faFileCodeO,
				125/\texttt{MUSTThesisS.sty}/\faFileCodeO,
				175/\texttt{readme.tex}/\faFileCodeO,
				225/\texttt{ref.tex}/\faFileCodeO
			}
			{
				\node(file\a)[left=0pt of \ff.south west,yshift=-\a * .1mm,label=left:\color{gray}\texttt{{\footnotesize \b}},black!70]{ \c };
				\draw[line width=1pt,rounded corners=2.25pt,gray](\ff) |-(file\a);
			}




		% reference
		\ff{folder5}
		\node(\ff)[left=0pt of project.south west, yshift=-7cm,label={left:\color{blue}\texttt{\textbf{reference}}}]{{\Large \faFolderOpenO}};
		\draw[line width=1pt,rounded corners=1pt,gray](project) |-(\ff);

		\foreach \a /\b /\c  in {
				25/\texttt{ch-ref.bib}/\faFileTextO
			}
			{
				\node(file\a)[left=0pt of \ff.south west,yshift=-\a * .1mm,label=left:\color{gray}\texttt{{\footnotesize \b}},black!70]{ \c};
				\draw[line width=1pt,rounded corners=2.25pt,gray](\ff) |-(file\a);
			}



		% Tex 主文档
		\ff{folder5}
		\node(\ff)[
			right=0pt of project.south east, yshift=-6.4cm,
			label={[draw=Cyan,line width=1.2pt,text width=2.6cm]right:\color{black!60}\texttt{\textbf{MUSTThesis.tex}}},
			black!70
		]{ \faFileTextO};

		\node(df)[right=3cm of \ff.east]{\#~ \LaTeX 主文档};

		\draw[line width=1pt,rounded corners=1pt,gray](project) |-(\ff);


		% lic许可文件
		\ff{folder6}
		\node(\ff)[
			right=0pt of project.south east, yshift=-7.1cm,
			label={[draw=Cyan,line width=1.2pt,text width=2.6cm]right:\color{gray}\texttt{lppl-1-3c.lic}},
			black!70]
		{ \faFileTextO};

		\node(df)[right=3cm of \ff.east]{\#~许可文件};
		\draw[line width=1pt,rounded corners=1pt,gray](project) |-(\ff);


		% iitool.exe
		\ff{folder7}
		\node(\ff)[
			right=0pt of project.south east, yshift=-7.8cm,
			label={[draw=Cyan,line width=1.2pt,text width=2.6cm]right:\color{gray}\texttt{iitool.exe}},
			black!70]
		{ \faFileTextO};
		\node(df)[right=3cm of \ff.east]{\#~自动化编译工具};
		\draw[line width=1pt,rounded corners=1pt,gray](project) |-(\ff);



	\end{tikzpicture}
\end{center}

\vfill







\restoregeometry
\newpage